\section{Introducción}
\label{sec:introduccion}

El presente informe desarrolla la Actividad 6 de Sistemas de Radiocomunicación \cite{consigna}, dedicada a los medios de comunicación analógicos y digitales. La actividad tiene tres partes:
\begin{itemize}
    \item \underline{Parte 1:} mediciones con el analizador de espectro sobre señales recibidas de radiodifusión en AM, en FM, de televisión analógica y de televisión digital terrestre (TDA, norma ISDB-Tb). Para cada servicio se identifica la emisora o el canal con el Cuadro de Atribución de Bandas de Frecuencias de la República Argentina (CABFRA) \cite{cabfra}, se estima el ancho de banda ocupado y se interpreta el espectro según la modulación utilizada.
    \item \underline{Parte 2:} interpretación de los niveles expresados en dBm, dBµV y dBc, conversión entre ellos para sistemas de $50$ y $75 \, [\Omega]$, y cálculo del índice de modulación de una señal AM a partir de los niveles de portadora y de banda lateral.
    \item \underline{Parte 3:} investigación de las tecnologías que se usan para distribuir televisión y datos por cable y por fibra óptica: DVB-C, DOCSIS, Video RF Overlay en redes FTTH, ISDB-Tb e ISDB-Tb sobre cable.
\end{itemize}

La base teórica proviene de los videos de la asignatura \cite{videos_medios} y de la bibliografía de comunicaciones \cite{tomasi, haykin}. Los temas de la Parte 3 se apoyan en las normas de cada sistema y en documentación de fabricantes, que se citan en cada caso.

\subsection{Objetivos}
\begin{itemize}
    \item Utilizar el analizador de espectro para identificar y medir señales de radiofrecuencia, registrando en cada caso su configuración.
    \item Reconocer las características espectrales de las emisiones de AM, de FM y de televisión analógica y digital, y relacionarlas con la modulación y con la canalización de cada servicio.
    \item Interpretar los niveles de señal expresados en dBm, dBµV y dBc, y convertirlos entre sí teniendo en cuenta la impedancia.
    \item Describir las principales tecnologías de distribución de televisión y datos por redes de cable y de fibra óptica, con mayor detalle en el Video RF Overlay sobre FTTH, que complementa el diseño de la red FTTH del Trabajo Práctico 4 \cite{tp4}.
\end{itemize}
